\begin{claim}{C.7c}{Proved}
Let $p$ be any prime, $K/\Q_p$ finite, $F_K=K(\mu_\infty)$, and
$H_K=G_{F_K}$. For the tripod $U=\mathbf P^1_K\setminus\{0,1,\infty\}$,
let $\Pi$ be its full geometric group and $P=\Pi^{(p)}$. Then
\[
 \operatorname{im}(\Pi^{s_t(H_K)}\longrightarrow P)=1
       \quad(t\in U(K)),\qquad
 \operatorname{im}(\Pi^{s_{\overrightarrow{01}}(H_K)}
                    \longrightarrow P)=I_0^{(p)}.
\]
The assertions concern images of the full fixed groups. In particular they
do not assert that all invariants of the maximal pro-$p$ quotient lift.
\end{claim}
\begin{proof}
We use two elementary refinements of the Magnus argument in C.7a. First,
for a free pro-$p$ group of any finite rank, its Magnus homomorphism into
$1+(X_1,\ldots,X_r)\Z_p\langle\!\langle X_1,\ldots,X_r\rangle\!\rangle$
is injective by the same finite-$p$ group-algebra argument given there.
If an automorphism acts identically in degree one and its first
homogeneous derivation is $\operatorname{ad}(L)$ for a nonzero linear
polynomial $L$, a fixed element with zero abelianization is trivial.
Indeed its first nonzero homogeneous Magnus term $W$ would commute with
$L$. A linear change of variables over $\Q_p$ and the word argument in
C.7a identify the polynomial centralizer as $\Q_p[L]$. Thus $W$ would be
a nonzero scalar multiple of $L^d$. But commutative specialization of the
entire Magnus expansion of an element with zero abelianization is $1$;
its degree-$d$ term cannot specialize to that nonzero polynomial.

Second, if finitely many continuous homomorphisms
$\lambda_1,\ldots,\lambda_r,\lambda:H_K\to\Q_p$ are linearly
independent, there is $\sigma\in H_K$ with
$\lambda_i(\sigma)=0$ for all $i$ and $\lambda(\sigma)\ne0$.
Their joint image is a compact $\Z_p$-submodule of $\Q_p^{r+1}$ of full
rank, hence contains a nonzero multiple of the last coordinate vector.
This proves the assertion. It applies to independent local Kummer
classes by the injectivity (C.7b.1).

For an ordinary tripod point $t$, the first derivation of its genuine
point action on $H_K$ is
\[
 \operatorname{ad}\bigl(\kappa(t)(\sigma)X+
                         \kappa(1-t)(\sigma)Y\bigr).           \tag{C.7c.1}
\]
The tangential action contributes no degree-one derivation because its
conjugating word lies in the commutator, as in C.7a. If at least one of
$t,1-t$ is not a root of unity, local Kummer injectivity supplies a
$\sigma$ for which the linear polynomial in (C.7c.1) is nonzero. C.7b
supplies zero full abelianization, so the preceding Magnus observation
proves the ordinary-point assertion in this case. Both $t$ and $1-t$
are roots of unity exactly when $t^2-t+1=0$: embed their number field in
$\C$ and intersect the two unit circles centred at $0$ and $1$.

We handle this exceptional case by a rational cyclic cover. We record the
first-derivation calculation used here and below. On a punctured projective
line whose finite cusps are $c_i$, compare an ordinary base point $x$ to
the unit tangent at a cusp $c_i$. The path cocycle's abelian coordinates
are given by the values of the invertible functions $z-c_j$: for $j\ne i$
the coordinate is
$\kappa(c_i-c_j)-\kappa(x-c_j)$, and for $j=i$ it is
$-\kappa(x-c_i)$. This follows by applying the Kummer coverings of these
functions to the path. For a tangential base at $0$, replace its value
of $z$ by its leading coefficient and retain the ordinary values of the
other functions. Thus, if all nonzero cusp differences have Kummer value
zero on a particular $\sigma$, the unit-tangent action at $0$ has no
first derivation. An ordinary point then contributes precisely the inner
derivation given by its own linear path term. The conjugation convention
may reverse all these signs simultaneously, which leaves every kernel
argument unchanged.

Suppose first $p\ge5$. Choose $c\in\mu_6$ with $c^p=t$, a primitive
$p$th root $\zeta$, and the lift $\alpha=c\zeta$ of $t$ on $z^p=x$.
Put $M_0=\Q_p(\mu_6,\mu_p)$. There is an automorphism $\tau$ of $M_0$
fixing $\mu_6$ and inverting $\mu_p$: the first extension is unramified,
whereas $\Q_p(\mu_p)/\Q_p$ is totally ramified. Every rational Kummer
class of a difference of two $p$th roots of unity is $\tau$-fixed, since
\[
 \frac{\zeta_i^{-1}-\zeta_j^{-1}}{\zeta_i-\zeta_j}
       =-\frac1{\zeta_i\zeta_j}
\]
is a root of unity. Differences from the cusp $0$ are roots of unity
and have zero class. On the other hand $u=\kappa(1-c\zeta)$ is not
$\tau$-fixed. The ratio
\[
 \frac{1-c\zeta}{1-c\zeta^{-1}}
\]
is not a root of unity: in a complex embedding with
$c=\exp(\pm\pi i/3)$ and $\zeta=\exp(2\pi i/p)$, the squared absolute
values of numerator and denominator differ by a nonzero multiple of
$\sin(\pi/3)\sin(2\pi/p)$. Local rational Kummer theory therefore makes
their difference nonzero.

For $p=2$ or $3$, use $z^2=x$, $\alpha^2=t$, and
$M_0=\Q_p(\mu_{12})$. There is an automorphism $\tau$ fixing $\mu_6$
and sending $\alpha$ to $-\alpha$. At $p=2$ this is the ramified quadratic
automorphism on $\mu_4$, independent of the unramified $\mu_3$; at $p=3$
it is the unramified quadratic automorphism on $\mu_4$, independent of the
ramified $\mu_3$. The only nonzero cusp-difference class on the quadratic
cover is generated by $\kappa(2)$, which is $\tau$-fixed. But
$\kappa(1-\alpha)$ is nonzero and $\tau$ sends it to its negative:
$(1-\alpha)(1+\alpha)=1-t$ is a root of unity, whereas $1-\alpha$ is not,
as its complex absolute value shows.

In either case the selected ordinary-point path class lies outside the
span of the cusp-difference classes. Pass to $M=KM_0$ if necessary;
restriction on rational Kummer cohomology remains injective, since
corestriction after restriction multiplies by the finite degree. Then
(C.7b.1) keeps these classes independent on $H_K$, because $M_0$ is
cyclotomic and $M(\mu_\infty)=F_K$. Choose a $\Q_p$-basis among the cusp-difference classes, discarding
any dependent classes. The selected point-path class is independent of
that basis, so the compact-lattice observation supplies a $\sigma\in H_K$
annihilating the basis, hence every cusp-difference class, but not the
selected point-path class. The
first-derivation calculation gives a nonzero inner derivation on the
free pro-$p$ group of the covering punctured projective line.
A full fixed $g$ has zero full winding by C.7b, so lies in this cyclic
covering subgroup. Its lift is fixed at $\alpha$, and has zero full
abelianization by C.7b applied to the covering curve. The Magnus
observation makes its pro-$p$ image trivial, hence makes the original
image in $P$ trivial. This finishes the ordinary-point assertion.

For the tangential assertion we need a homogeneous derivation lemma.
Let $A$ be the free associative algebra over a characteristic-zero field
on finitely many variables, partitioned into at least two nonempty
clusters. Write $S_C$ for each cluster's sum and $S=\sum_C S_C$.
Suppose
\[
 D=c\operatorname{ad}(S)+\sum_C D_C,\qquad c\ne0,              \tag{C.7c.2}
\]
where $D_C$ is a homogeneous degree-one derivation involving only variables
of $C$, zero on other clusters, and satisfies $D_C(S_C)=0$. Then
\[
 \ker(D)=\mathbf k[S].                                        \tag{C.7c.3}
\]
Here is a full block-length proof. A word is uniquely a sequence of
maximal nonempty blocks in the individual cluster algebras; consecutive
blocks have different clusters. The internal derivations do not increase
this number of blocks. Let $W$ be homogeneous of total degree $d$, with
maximal block length $b$. If $D(W)=0$, the block-length-$b+1$ part of
$[S,W]$ vanishes. For each alternating cluster sequence
$\mathbf i=(i_1,\ldots,i_b)$, let $W_{\mathbf i}$ be the corresponding
tensor component, in the tensor product of the cluster augmentation
ideals. The vanishing equations are
\[
 S_j\otimes W_{i_1,\ldots,i_b}
   =W_{j,i_1,\ldots,i_{b-1}}\otimes S_{i_b}
       \qquad(j\ne i_1).                                     \tag{C.7c.4}
\]
The tensor equation first forces the last factor of $W_{\mathbf i}$
to lie in the one-dimensional span of $S_{i_b}$. Repeating the shift in
(C.7c.4) forces every factor into the corresponding span. Thus
$W_{\mathbf i}=\lambda_{\mathbf i}S_{i_1}\otimes\cdots\otimes S_{i_b}$.
The shifts connect all alternating cluster sequences: with two clusters
there are two alternating sequences and the shift interchanges them;
with three or more, an extra intermediate cluster permits any prescribed
sequence to be prepended. Hence all scalars are equal. The maximal-block
part is exactly a scalar multiple of the maximal-block part of $S^b$.
Since each $S_i$ has degree one, a nonzero component forces $d=b$.
Subtract that scalar multiple of $S^d$, which lies in $\ker D$, from $W$.
Any nonzero remainder has total degree $d$ and maximal block length
strictly less than $d$, which the same argument excludes. This proves
(C.7c.3), including constants.

Apply this to
\[
 W_n=\mathbf P^1_K\setminus(\mu_{p^n}\cup\{\infty\}),
       \qquad \hbox{ordinary base point }0.
\]
Extend $K$ to $M=K(\mu_{p^n})$. Choose a uniformizer $\pi_M$ and units
whose rational Kummer classes form a basis of the unit subspace of
$H^1(G_M,\Q_p(1))$. By (C.7b.1) and the compact-lattice observation,
there exists $\sigma\in H_K$ whose Kummer evaluations kill every unit
class and are nonzero on $\pi_M$. They are therefore a common nonzero
multiple of the normalized valuation $v_M$ on $M^\times$.
Use free pro-$p$ cusp generators indexed by $\zeta_i\in\mu_{p^n}$.
The preceding path calculation, now starting at the ordinary point $0$,
gives the first derivation, up to a common nonzero scalar and sign,
\[
 D(X_i)=\sum_{j\ne i}v_M(\zeta_i-\zeta_j)[X_j,X_i].           \tag{C.7c.5}
\]
The starting values $-\zeta_j$ and each terminal tangent coefficient
are roots of unity, so there are no other terms.

Index roots by integers modulo $p^n$ and partition them according to their
index modulo $p$. Differences in distinct clusters have the same minimal
positive valuation
\[
 c=\frac{e(M/\Q_p)}{p^{n-1}(p-1)}\ne0.
\]
Indeed $\Phi_{p^r}(1+T)$ is Eisenstein over $\Q_p$, and the norm
of $1-\zeta_{p^r}$ is $p$, giving valuation
$e(M/\Q_p)/(p^{r-1}(p-1))$. For a difference of roots indexed by
$i,j$, take $r=n-v_p(i-j)$; roots from distinct listed clusters have
$r=n$. Subtract this common value from every coefficient in (C.7c.5).
The remaining terms involve only pairs within the same cluster.
Their coefficients are symmetric in $i,j$, so their derivation preserves
the cluster sum. Thus (C.7c.5) has exactly the form (C.7c.2).
If a full fixed element on $W_n$ had a nontrivial maximal pro-$p$ image,
its first nonzero Magnus term would lie in $\ker D$ and hence be a
nonzero multiple of $S^d$. This contradicts its zero full abelianization
from C.7b, by commutative specialization. We have proved
\[
 \operatorname{im}\bigl(\pi_1(W_{n,\bar K})^{s_0(H_K)}
                   \to\pi_1(W_{n,\bar K})^{(p)}\bigr)=1.      \tag{C.7c.6}
\]

Finally let $g$ be a full fixed element for the unit tangent on $U$.
Subtract the cusp-$0$ element with the same winding, obtaining $h$ with
zero $a$-coordinate. For every $n$ it lifts to the unit tangent at $0$
on the power cover $z^{p^n}=x$. More explicitly, choose compatible
Puiseux roots $x^{1/p^n}$ with leading coefficient $1$ and
$(x^{1/p^{n+1}})^p=x^{1/p^n}$. Sending $z$ to $x^{1/p^n}$ defines the
lifted unit tangent; the Galois action is on coefficients and fixes these
formal parameters, so the original section restricts to that lifted
section on $H_K$. Fill in $0$. Formula (C.7c.6) kills its
image in the maximal pro-$p$ group of $W_n$. The kernel of filling in
that cusp is the closed normal subgroup generated by its cusp inertia;
this is immediate from the free geometric presentation. Its generator
maps to $a^{p^n}$ in the original $P$. Consequently
\[
 \overline h\in\bigcap_{n\ge1}
       \overline{\langle\!\langle a^{p^n}\rangle\!\rangle}_P=1.
\]
The last intersection is trivial because, in every finite $p$-group
quotient of $P$, the image of $a^{p^n}$ is trivial for all sufficiently
large $n$. We obtain the asserted cusp upper bound; the actual full cusp
inertia is fixed by $H_K$ and gives equality of the image.
\end{proof}
